\phantomsection
\addcontentsline{toc}{section}{Glossary of terms}
\label{sec:glossary}
\section*{Glossary of terms}

A reader does not need to know every term below to follow the document; this list is here to look one
up when it appears. Concepts are explained in plain language; acronyms are expanded.

\begin{description}[leftmargin=2.2em,style=nextline,font=\normalfont\bfseries]
  \item[ATT / GATT] The Bluetooth Low Energy layers that define ``characteristics'' (named values a
        device exposes) and the services that group them.
  \item[BLE (Bluetooth Low Energy)] A short-range, low-power wireless protocol; the collar's only
        radio.
  \item[BMS] Battery Management System: the protection circuit built into the LiPo cell.
  \item[Confidence-gated cascade] The collar's two-model scheme: a motion model runs on every window,
        and a second audio model is consulted only when the motion model is unsure.
  \item[CMSIS-NN] ARM's hand-optimised neural-network routines for Cortex-M microcontrollers.
  \item[CRC] Cyclic Redundancy Check: a checksum used to verify an over-the-air model transfer.
  \item[ESP-IDF] Espressif's official development framework (SDK) for the ESP32 family; the station's
        software stack.
  \item[Firebase] Google's hosted backend platform: user sign-in, a database, file storage, and
        serverless functions, used here as the cloud tier.
  \item[Firmware] The program that runs on a microcontroller (the collar and the station each run
        their own).
  \item[IMU (Inertial Measurement Unit)] The collar's six-axis motion sensor: a three-axis
        accelerometer plus a three-axis gyroscope.
  \item[Inference] Running a trained model on new input to get a prediction (here, classifying a
        behaviour on the collar).
  \item[int8 quantisation] Storing a model's numbers as 8-bit integers instead of floating point,
        which shrinks it and speeds it up enough to run on a microcontroller.
  \item[JWT] JSON Web Token: the signed token format of a Firebase login.
  \item[MTU] Maximum Transmission Unit: the negotiated size of a BLE data packet (247 bytes here).
  \item[NCS (nRF Connect SDK)] Nordic's distribution of the Zephyr operating system, used to build
        the collar firmware.
  \item[NimBLE] The lightweight Bluetooth stack used in the station firmware.
  \item[NVS] Non-Volatile Storage: the station's small flash key/value store for its credentials.
  \item[OTA (Over The Air)] Delivering a new trained model to the collar wirelessly over Bluetooth,
        with no cable or reflash.
  \item[PDM] Pulse-Density Modulation: the raw output format of the collar's MEMS microphone.
  \item[PHY] The Bluetooth physical layer; the collar uses the 2\,Mbit variant for speed.
  \item[RSSI] Received Signal Strength Indication: how strong the collar's signal is at the station,
        used as a ``close enough to record'' gate.
  \item[RTDB (Realtime Database)] Firebase's live, JSON-tree database, holding each cat's current
        state and history.
  \item[SoC (System on Chip)] A single chip combining processor, memory and radio , the nRF52840 and
        ESP32-S3 modules.
  \item[Tensor arena] A fixed block of RAM the inference engine uses as scratch space for one model.
  \item[TFLM (TensorFlow Lite Micro)] The on-device engine that runs the trained models on the collar.
  \item[UF2] The drag-and-drop file format used to flash the collar over USB.
  \item[Web Serial] A browser feature that lets a web page talk to a USB device, used to pair the
        station.
  \item[Zephyr] A real-time operating system for microcontrollers; the base of the collar firmware.
  \item[$\mu$-law] An 8-bit audio compression scheme (G.711) that halves the audio sent over
        Bluetooth.
\end{description}
\clearpage
